\chapter{Dataset Description}
\label{ch:dataset}

\section{Origin of the data}
All training and evaluation data in this project derive from MailEx \citep{srivastava2023mailex}. The emails themselves come from the Enron corpus \citep{klimt2004enron}: mailbox file names such as \code{bass-e\_inbox114} and the corporate disclaimers in many messages make this visible in every record. MailEx adds event annotations: each email is labelled with zero or more events of a typed taxonomy, and each event with a trigger and argument spans such as \lab{Action Description}, \lab{Action Members} or \lab{Meeting Date}.

MailEx was chosen for three reasons. It is, to our knowledge, the only public corpus with argument-level annotation of requests and deliveries in email threads; its arguments correspond closely to the information a task list needs (what, who, when); and its thread structure allows experiments with conversation context and thread-disjoint splits. Its limitations for this purpose are discussed in \cref{sec:datalimits}.

\section{Original format and unified representation}
The MailEx release stores threads as JSON files with token lists and per-event BIO label sequences whose labels concatenate event type and argument role, for example \code{Request\_Action:B-Action Description}. The conversion script \code{scripts/convert\_mailex.py} turns each email (one turn of a thread) into a unified JSONL record with the fields listed in \cref{lst:unified}. Character offsets for every token are recovered and checked, so that any later span can be traced back to the exact characters of the original email.

\begin{lstlisting}[language=json, caption={Fields of a unified MailEx record (values abbreviated from \code{data/unified/train.jsonl}).}, label={lst:unified}]
{"id": "mailex:train:bass-e_inbox114:turn_0000",
 "source_split": "train", "thread_id": "bass-e_inbox114", "turn_index": 0,
 "text": "For some reason your e-mail to Shanna and I was quarantined ...",
 "tokens": ["For", "some", "reason", ...], "token_offsets": [[0,3], [4,8], ...],
 "sender": "\"Bass, Eric\" <Eric.Bass@ENRON.com>", "recipients": ["<...>"],
 "timestamp": null, "subject": "Your E-mail",
 "thread_context": {"turn_count": 2, "next_record_id": "...turn_0001", ...},
 "original_events": {"Request_Data": {"labels": [["O", ..., "Request_Data:Context: B-Data idString", ...]],
                                      "triggers": [...]}},
 "text2task_annotations": null}
\end{lstlisting}

Two properties of the records shaped later design decisions. First, \code{timestamp} is \code{null} for every record: MailEx does not preserve sending times, so relative dates in the corpus cannot be resolved to calendar days (\cref{sec:temporal}). Second, sender and recipient fields are raw header strings, which makes them a low-trust source of participant identities.

\section{Size and splits}
\Cref{tab:mailexsplits} gives the split sizes recomputed directly from the unified files. They match the figures published by the MailEx authors exactly (3,117/414/405 emails; 1,200/150/150 threads), and the number of emails carrying only the null event label \lab{O} (636/70/70) matches their count of non-event emails. A further 272, 25 and 30 emails have empty text after conversion; these are kept in the unified files but skipped by every model-ready builder.

\begin{table}[h]
\centering
\caption{MailEx splits as converted (recomputed from \code{data/unified/*.jsonl}).}
\label{tab:mailexsplits}
\small
\begin{tabular}{@{}lrrrr@{}}
\toprule
Split & Threads & Emails & Emails without events & Empty text \\
\midrule
Train & 1,200 & 3,117 & 636 & 272 \\
Development & 150 & 414 & 70 & 25 \\
Test & 150 & 405 & 70 & 30 \\
\midrule
Total & 1,500 & 3,936 & 776 & 327 \\
\bottomrule
\end{tabular}
\end{table}

\Cref{fig:eventtypes} shows how often each event type occurs. Deliveries of information (\lab{Deliver\_Action\_Data}, \lab{Deliver\_Data}) dominate; explicit requests for actions (\lab{Request\_Action}) occur in 834 training emails, and meeting-related events are comparatively rare. Eleven distinct event-type names occur in the converted files (the paper groups them into ten types).

\begin{figure}[t]
\centering
\chart{mailex_event_types}
\caption[MailEx event types by split]{Number of emails containing at least one event of each type, per split. ``No event (O)'' counts emails whose only label is the null event. Recomputed from the unified files.}
\label{fig:eventtypes}
\end{figure}

\section{Label sets used in this project}
The project produced three label sets over these emails, each with a different origin and quality level (\cref{tab:labelsets}). Their construction is explained in \cref{ch:preprocessing}; here we describe what they contain.

\begin{table}[h]
\centering
\caption{The three label sets. Major-error rates come from AI-assisted inspections of 30 to 40 random emails per set and are indicative only.}
\label{tab:labelsets}
\footnotesize
\begin{tabularx}{\linewidth}{@{}p{26mm}p{22mm}p{26mm}Lp{18mm}@{}}
\toprule
Label set & Emails (train/dev/test) & Entity types & Origin & Major-error rate \\
\midrule
SILVER-v1 & 450 / 60 / 90 & TASK, PERSON, TEMPORAL, ARTIFACT, GROUP, EVENT; 4 relation types & Automatic and LLM annotation of 600 selected candidates, schema validation & 22 of 30 (73\%) \\
SILVER-v2 & 449 / -- / -- & TASK, PERSON, TEMPORAL & Repair of v1: audited records, validated LLM outputs, conservative filtering & 10 of 30 (33\%) \\
MailEx-derived & 2,845 / 389 / not built & TASK, PERSON, TEMPORAL & Deterministic mapping of MailEx human arguments & about 9\% (pooled) \\
\bottomrule
\end{tabularx}
\end{table}

SILVER-v2 has no validation split of its own: models trained on it are validated on the 60 SILVER-v1 validation emails transformed to the three-class view. The MailEx-derived test split was deliberately never built, so no model has been scored on held-out test data in that lineage.

\subsection{Entity distributions}
\Cref{fig:entitydist} compares entity counts across the label sets. Two observations matter for modelling. TASK is the most frequent type in every set, but its spans are long: in the MailEx-derived training set the 1,664 TASK spans cover 17,028 WordPiece tokens (about 10 tokens per span), whereas PERSON spans average about 1.5 tokens. And the six-class SILVER-v1 set contains very few GROUP (35) and EVENT (4) training spans, too few to learn from.

\begin{figure}[t]
\centering
\chart{entity_distribution}
\caption[Entity distributions across label sets]{Entity span counts in the three label sets. SILVER-v1 is the six-class set; SILVER-v2 and the MailEx-derived sets use three classes.}
\label{fig:entitydist}
\end{figure}

\subsection{Relation distribution}
Relations are available only in the SILVER-v1 lineage. The relation dataset enumerates every type-valid task--person and task--temporal pair in an email as a candidate and labels it with the annotated relation, or \lab{NONE}. \Cref{fig:reldist} shows the result. Negatives outnumber positives about four to one, \lab{DUE\_ON} is extremely rare (33 training pairs, 10 validation pairs), and \lab{DEPENDS\_ON} has only 4 training examples, all linking a TASK to an EVENT, which the three-class entity model cannot produce. This is why dependencies are handled by rules (\cref{ch:rules}).

\begin{figure}[t]
\centering
\chart{relation_distribution}
\caption[Relation label distribution]{Candidate pairs per relation label (log scale) over all type-valid TASK$\to$PERSON and TASK$\to$TEMPORAL pairs. For training, \lab{NONE} pairs were subsampled to 1,144 (1:1 against positives).}
\label{fig:reldist}
\end{figure}

\section{Representative annotated examples}
The following MailEx-derived training emails illustrate both the scheme and its noise. Spans are shown as \lab{[TYPE text]}.

\begin{enumerate}
  \item \emph{``I just chatted with [PERSON Andrea]. We[TASK 'll cover this] [TEMPORAL tomorrow].''} A commitment: the \lab{TASK} span begins at the contraction \emph{'ll} because the MailEx \lab{Action Description} argument does. \Cref{fig:bio} follows this example token by token.
  \item \emph{``[PERSON Louis], The dash looks good. \ldots{} Please [TASK try to push this, with Chip's help], [TEMPORAL tomorrow].''} A request with an addressee. \emph{Chip} is not labelled, because MailEx marks people only when they are event arguments.
  \item \emph{``[PERSON Shelley] [TASK will only be available] [TEMPORAL tomorrow] [TEMPORAL from 3:00 until 5:00].''} A statement of availability labelled as a TASK by the commitment rule: an example of label noise.
  \item \emph{``[PERSON Steve] -- [PERSON Christi] left you [TEMPORAL off this list]. Please [TASK include your thoughts on this worksheet as well for Eastern gas].''} The TASK is correct; \emph{off this list} is a MailEx date argument that is not a date at all.
\end{enumerate}

\section{Class imbalance}
\label{sec:imbalance}
At the token level the label distribution is extremely skewed (\cref{fig:bioimbalance}). Of 245,834 labelled WordPiece tokens in the MailEx-derived training set, 91.0\% are \lab{O}; each \lab{B-} label accounts for under 0.7\%. A model that predicts \lab{O} everywhere reaches about 91\% token accuracy while finding no entity at all, which is exactly what the first SILVER model did (token accuracy 0.73 on its validation set, entity F1 0.008; \cref{ch:evolution}). The imbalance is even stronger for relations of the rarest class, where \lab{DUE\_ON} pairs are outnumbered by \lab{NONE} pairs more than a hundredfold in the full candidate set.

\begin{figure}[t]
\centering
\chart[0.72\linewidth]{bio_imbalance}
\caption[BIO label imbalance]{WordPiece-level BIO label counts in the MailEx-derived training set (log scale), recomputed from \code{data/model\_ready/entities\_mailex\_v1/train.jsonl}.}
\label{fig:bioimbalance}
\end{figure}

\section{Limitations of the data}
\label{sec:datalimits}
\begin{itemize}
  \item \textbf{Different task definition.} MailEx annotates events, not tasks. Requests for \emph{data} (\emph{``Could you forward the article?''}) are annotated as \lab{Request\_Data} with the requested object as argument, so the mapping cannot recover them as TASK phrases.
  \item \textbf{Incompleteness.} Names and dates are labelled only when they are event arguments, and about one email in eight contains a clear request with no annotation at all. Measured precision is therefore understated and recall is unreliable.
  \item \textbf{Heuristic noise.} PERSON spans pass a name filter that still admits words such as \emph{director} or \emph{girl}; the commitment rule admits some predictions as tasks.
  \item \textbf{No timestamps.} Without sending times, relative dates in the corpus cannot be resolved; only 1.3\% of temporal spans resolve to a calendar value (\cref{sec:temporal}).
  \item \textbf{Almost no lifecycle or dependency examples.} Across 510 replayed emails there are about 17 completion cues, 8 cancellations, 11 reschedules and a single reassignment, and only 7 dependency annotations in the whole SILVER set. Reconciliation and dependency handling therefore cannot be trained or evaluated on corpus labels.
  \item \textbf{Duplicates.} Enron threads quote each other heavily. 39 of the 389 MailEx-derived development emails are exact or near copies of training emails (\cref{sec:leakage}).
  \item \textbf{Register.} Business email from 2001 differs from the short chat messages the application handles, a domain shift visible in \cref{ch:results}.
\end{itemize}
